% ============================================================================
%  TEMPLATE LaTeX — Rodapé Institucional UFMT
%  Disciplina : Algoritmos 1
%  Professora : Rafaela Souza Francisco
%  Curso      : Bacharelado em Ciência da Computação — UFMT, Cuiabá
% ============================================================================
\documentclass[12pt, a4paper]{article}

% --- Codificação e idioma ---------------------------------------------------
\usepackage[utf8]{inputenc}
\usepackage[T1]{fontenc}
% \usepackage[brazil]{babel}  % descomente se tiver texlive-lang-portuguese

% --- Fonte ------------------------------------------------------------------
% \usepackage{lmodern}        % descomente se disponível no seu sistema

% --- Layout de página -------------------------------------------------------
\usepackage[
  top=2.5cm,
  bottom=3.0cm,       % espaço extra para o rodapé
  left=2.5cm,
  right=2.5cm
]{geometry}

% --- Pacotes necessários ----------------------------------------------------
\usepackage{graphicx}          % inserir imagens
\usepackage[brazil]{babel}
\usepackage{fancyhdr}          % cabeçalho/rodapé customizado
\usepackage{xcolor}            % cores
\usepackage{hyperref}          % links clicáveis
\usepackage{tikz}              % desenhos (usado na linha decorativa)
\usepackage{amsmath, amssymb}
\usepackage{lastpage}          % referência à última página

\usepackage{booktabs}
\usepackage{enumitem}
\usepackage{array}
\usepackage{geometry} 


\usepackage{listings}
\usepackage{xcolor}
\renewcommand{\lstlistingname}{Código}

\lstdefinestyle{cstyle}{
  language=C,
  basicstyle=\ttfamily\small,
  keywordstyle=\bfseries\color{blue!70!black},
  commentstyle=\itshape\color{green!50!black},
  stringstyle=\color{orange!80!black},
  numbers=left,
  numberstyle=\tiny\color{gray},
  numbersep=8pt,
  frame=single,
  breaklines=true,
  tabsize=4,
  showstringspaces=false,
}

\lstset{
    basicstyle=\ttfamily\small,
    keywordstyle=\color{blue}\bfseries,
    commentstyle=\color{gray}\itshape,
    stringstyle=\color{red},
    numbers=left,
    numberstyle=\tiny\color{gray},
    numbersep=8pt,
    frame=single,
    breaklines=true,
    tabsize=4,
    showstringspaces=false
}

\usetikzlibrary{shapes.geometric, arrows}
\tikzstyle{iniciofim} = [ellipse, draw, minimum width=3cm, minimum height=1cm, text centered]
\tikzstyle{processo} = [rectangle, draw, minimum width=5cm, minimum height=1cm, text centered]
\tikzstyle{seta} = [->, thick]
\tikzstyle{decisao} = [diamond, minimum width=3.5cm, minimum height=2cm, text centered, draw=black]

% --- Cores institucionais UFMT ---------------------------------------------
\definecolor{ufmtAzul}{RGB}{0, 51, 102}       % azul escuro institucional
\definecolor{ufmtVerde}{RGB}{0, 105, 62}       % verde institucional
\definecolor{ufmtCinza}{RGB}{100, 100, 100}    % cinza para texto secundário

% --- Dados do documento (EDITE AQUI) ----------------------------------------
\newcommand{\NomeAluno}{Thalles Eduardo Silva Mota}             % ← Seu nome
\newcommand{\Disciplina}{Cálculo 1}
\newcommand{\Professora}{Prof. Saulo Alves de Araújo}
\newcommand{\Curso}{Bacharelado em Ciência da Computação}
\newcommand{\Instituicao}{UFMT — Universidade Federal de Mato Grosso}
\newcommand{\Semestre}{2026/1}                                  % ← Semestre atual
\renewcommand{\abstractname}{Resumo}

% Se tiver o logo em arquivo, defina o caminho abaixo:
% \newcommand{\LogoPath}{logo_ufmt.png}

% ============================================================================
%  CONFIGURAÇÃO DO RODAPÉ  (o coração do template)
% ============================================================================
\pagestyle{fancy}
\fancyhf{}                     % limpa cabeçalho e rodapé padrão

% --- Cabeçalho (opcional — descomente se quiser) ----------------------------
% \fancyhead[L]{\small\color{ufmtCinza}\Disciplina}
% \fancyhead[R]{\small\color{ufmtCinza}\NomeAluno}
% \renewcommand{\headrulewidth}{0.4pt}

% --- Rodapé -----------------------------------------------------------------
\renewcommand{\footrulewidth}{0pt}  % remove linha padrão (usaremos tikz)

\fancyfoot[C]{%
  \vspace{2pt}%
  % ── Linha decorativa com gradiente ──────────────────────────────────
  \begin{tikzpicture}[remember picture, overlay]
    \draw[
      line width=1.2pt,
      ufmtAzul
    ] ([yshift=8pt]current page.south west |- 0,0) ++(2.5cm,0)
      -- ++(\dimexpr\textwidth\relax, 0);
    % Pequeno detalhe em verde (identidade visual)
    \fill[ufmtVerde]
      ([yshift=6.5pt]current page.south west |- 0,0) ++(2.5cm,0)
      rectangle ++(4cm, 1.8pt);
  \end{tikzpicture}%
  %
  % ── Conteúdo do rodapé ──────────────────────────────────────────────
  \begin{minipage}[t]{0.12\textwidth}
    \vspace{-2pt}
    % =====================================================================
    % LOGO: Descomente a linha abaixo e comente o \fbox quando tiver o logo
    % \includegraphics[height=1.1cm]{\LogoPath}
    % =====================================================================
    % Placeholder visual enquanto não tiver o .png:
    \includegraphics[height=1.1cm]{images/UFMT - logo texto.png}%
  \end{minipage}%
  \hfill
  \begin{minipage}[t]{0.60\textwidth}
    \centering
    \vspace{-2pt}
    {\footnotesize\color{ufmtAzul}\textbf{\Instituicao}}\\[1pt]
    {\scriptsize\color{ufmtCinza}%
      \Curso\ \textbullet\ \Disciplina\ — \Semestre}\\[1pt]
    {\scriptsize\color{ufmtCinza}\Professora}
  \end{minipage}%
  \hfill
  \begin{minipage}[t]{0.15\textwidth}
    \raggedleft
    \vspace{-2pt}
    {\footnotesize\color{ufmtAzul}\textbf{\thepage}\,/\,\pageref{LastPage}}
  \end{minipage}%
}

% Estilo da primeira página (pode ser diferente)
\fancypagestyle{plain}{%
  \fancyhf{}%
  \renewcommand{\headrulewidth}{0pt}%
  \fancyfoot[C]{%
    \vspace{2pt}%
    \begin{tikzpicture}[remember picture, overlay]
      \draw[line width=1.2pt, ufmtAzul]
        ([yshift=8pt]current page.south west |- 0,0) ++(2.5cm,0)
        -- ++(\dimexpr\textwidth\relax, 0);
      \fill[ufmtVerde]
        ([yshift=6.5pt]current page.south west |- 0,0) ++(2.5cm,0)
        rectangle ++(4cm, 1.8pt);
    \end{tikzpicture}%
    \begin{minipage}[t]{0.12\textwidth}
      \vspace{-2pt}
      \fbox{\footnotesize\textbf{\color{ufmtAzul}UFMT}}%
    \end{minipage}%
    \hfill
    \begin{minipage}[t]{0.60\textwidth}
      \centering\vspace{-2pt}
      {\footnotesize\color{ufmtAzul}\textbf{\Instituicao}}\\[1pt]
      {\scriptsize\color{ufmtCinza}%
        \Curso\ \textbullet\ \Disciplina\ — \Semestre}\\[1pt]
      {\scriptsize\color{ufmtCinza}\Professora}
    \end{minipage}%
    \hfill
    \begin{minipage}[t]{0.15\textwidth}
      \raggedleft\vspace{-2pt}
      {\footnotesize\color{ufmtAzul}\textbf{\thepage}\,/\,\pageref{LastPage}}
    \end{minipage}%
  }%
}

% ============================================================================
%  DOCUMENTO DE EXEMPLO
% ============================================================================
\begin{document}

\begin{center}
  {\Large\bfseries\color{ufmtAzul} Lista 1: Conjuntos e Equações}\\[6pt]
  {\large \NomeAluno}\\[3pt]
  {\normalsize\color{ufmtCinza} \today}
\end{center}

\vspace{0.5cm}

\section*{Exercício 1 - Resolução}
 
\begin{enumerate}[label=(\alph*)]
 
    \item $5 \in \mathbb{N}$: é um inteiro não negativo, portanto pertence
    a $\mathbb{N}$, $\mathbb{Z}$ e $\mathbb{Q}$.
 
    \item $-3 \notin \mathbb{N}$: é um inteiro negativo, logo está em
    $\mathbb{Z}$ e $\mathbb{Q}$, mas \textbf{não} em $\mathbb{N}$.
 
    \item $\dfrac{2}{7} \notin \mathbb{Z}$: é uma fração com
    $\operatorname{mdc}(2,7) = 1$, portanto é racional mas não inteiro.
 
    \item $0 \in \mathbb{N}$: pela convenção moderna,
    $\mathbb{N} = \{0, 1, 2, \ldots\}$, logo $0 \in \mathbb{N} \subset
    \mathbb{Z} \subset \mathbb{Q}$.
 
    \item $-10 \in \mathbb{Z} \setminus \mathbb{N}$: inteiro negativo,
    mesma situação de $-3$.
 
    \item $-1/3 \notin \mathbb{Z}$: fração negativa, pertence apenas
    a $\mathbb{Q}$.
 
\end{enumerate}

\section*{Exercício 2 — Resolução}

\begin{enumerate}[label=(\alph*)]
 
    \item $-2 \in \mathbb{N}$ \quad \textbf{(F)}\\
    
    $\mathbb{N} = \{0, 1, 2, \ldots\}$ contém apenas inteiros não negativos.
    Como $-2 < 0$, temos $-2 \notin \mathbb{N}$.
    
 
    \item $\dfrac{1}{2} \in \mathbb{Q}$ \quad \textbf{(V).}\\
    
    Todo número da forma $\dfrac{p}{q}$ com $p, q \in \mathbb{Z}$ e $q \neq 0$
    é racional. Como $\dfrac{1}{2}$ satisfaz isso, $\dfrac{1}{2} \in \mathbb{Q}$.
    
 
    \item $3 \in \mathbb{Z}$ \quad \textbf{(V)}\\
    
    $\mathbb{Z} = \{\ldots, -2, -1, 0, 1, 2, 3, \ldots\}$, portanto $3 \in \mathbb{Z}$.
    Além disso, como $\mathbb{N} \subset \mathbb{Z}$, isso é consequência direta
    de $3 \in \mathbb{N}$.
    
 
    \item $0 \in \mathbb{N}$ \quad \textbf{(V)}\\
    
    Pela convenção moderna, $\mathbb{N} = \{0, 1, 2, \ldots\}$,
    logo $0 \in \mathbb{N}$.
\pagebreak

\section*{Exercício 3 — Resolução}
 
\begin{enumerate}[label=(\alph*)]
 
    \item $\{x \in \mathbb{R} \mid x > 2\}$
 
    \textbf{Notação de intervalo:} $]2, +\infty[$
 
    \textit{Parêntese} em 2 pois $x > 2$ (extremo \textbf{excluído}).
 
    \begin{center}
    \begin{tikzpicture}
        \draw[->] (-0.5,0) -- (4,0) node[right] {$\mathbb{R}$};
        \draw (2, 0.08) -- (2, -0.08) node[below] {$2$};
        \draw[very thick, blue, ->] (2,0) -- (4,0);
        \draw[blue, fill=white] (2,0) circle (3pt); % círculo aberto = excluído
    \end{tikzpicture}
    \end{center}
 
    \item $\{x \in \mathbb{R} \mid -1 \leq x < 3\}$
 
    \textbf{Notação de intervalo:} $[-1, 3[$
 
    \textit{Colchete} em $-1$ pois $x \geq -1$ (extremo \textbf{incluído});
    \textit{parêntese} em $3$ pois $x < 3$ (extremo \textbf{excluído}).
 
    \begin{center}
    \begin{tikzpicture}
        \draw[->] (-2.5,0) -- (4,0) node[right] {$\mathbb{R}$};
        \draw (-1, 0.08) -- (-1, -0.08) node[below] {$-1$};
        \draw (3,  0.08) -- (3,  -0.08) node[below] {$3$};
        \draw[very thick, blue] (-1,0) -- (3,0);
        \draw[blue, fill=blue]  (-1,0) circle (3pt); % círculo fechado = incluído
        \draw[blue, fill=white] (3,0)  circle (3pt); % círculo aberto  = excluído
    \end{tikzpicture}
    \end{center}
 
    \item $\{x \in \mathbb{R} \mid x \leq 4\}$
 
    \textbf{Notação de intervalo:} $]-\infty, 4]$
 
    \textit{Colchete} em $4$ pois $x \leq 4$ (extremo \textbf{incluído}).
    O símbolo $-\infty$ nunca é incluído, usando-se sempre parêntese.
 
    \begin{center}
    \begin{tikzpicture}
        \draw[->] (-2,0) -- (5,0) node[right] {$\mathbb{R}$};
        \draw (4, 0.08) -- (4, -0.08) node[below] {$4$};
        \draw[very thick, blue, <-] (-2,0) -- (4,0);
        \draw[blue, fill=blue] (4,0) circle (3pt); % círculo fechado = incluído
    \end{tikzpicture}
    \end{center}
 
\end{enumerate}
\end{enumerate}

\section*{Exercício 4 — Resolução}

\begin{enumerate}[label=(\alph*)]

    \item $2x + 3 > 7$
    \[
        2x > 7 - 3 \implies 2x > 4 \implies x > 2
    \]
    \[
        \boxed{S = ]2, +\infty[}
    \]

    \item $3x - 5 \leq 4$
    \[
        3x \leq 4 + 5 \implies 3x \leq 9 \implies x \leq 3
    \]
    \[
        \boxed{S = ]-\infty, 3]}
    \]

    \item $-2x > 6$

    Dividindo por $-2$: \textbf{o sinal da desigualdade se inverte.}
    \[
        x < \frac{6}{-2} \implies x < -3
    \]
    \[
        \boxed{S = ]-\infty, -3[}
    \]

    \item $4 - x \geq 1$
    \[
        -x \geq 1 - 4 \implies -x \geq -3 \implies x \leq 3
    \]
    \[
        \boxed{S = ]-\infty, 3]}
    \]

\end{enumerate}
\pagebreak

\section*{Exercício 5 — Resolução}

\begin{enumerate}[label=(\alph*)]

    \item $|x| < 2$

    Inequação do tipo $|x| < k$ com $k > 0$:
    \[
        -2 < x < 2
    \]
    \[
        \boxed{S = ]-2,\ 2[}
    \]

    \item $|x| \leq 3$

    Inequação do tipo $|x| \leq k$:
    \[
        -3 \leq x \leq 3
    \]
    \[
        \boxed{S = [-3,\ 3]}
    \]

    \item $|x| > 1$

    Inequação do tipo $|x| > k$ com $k > 0$: gera \textbf{união} de dois intervalos.
    \[
        x < -1 \quad \text{ou} \quad x > 1
    \]
    \[
        \boxed{S = ]-\infty,\ -1[ \cup ]1,\ +\infty[}
    \]

    \item $|x| \geq 4$

    Inequação do tipo $|x| \geq k$:
    \[
        x \leq -4 \quad \text{ou} \quad x \geq 4
    \]
    \[
        \boxed{S = ]-\infty,\ -4\,] \cup [\,4,\ +\infty[}
    \]

\end{enumerate}

\section*{Exercício 6 --- Resolução}

\begin{enumerate}[label=(\alph*)]

    \item $|x - 1| < 2$
    \[
        -2 < x - 1 < 2 \implies -1 < x < 3
    \]
    \[
        \boxed{S = \left]-1,\ 3\right[}
    \]

    \item $|x + 3| \geq 5$
    \[
        x + 3 \leq -5 \implies x \leq -8
        \qquad \text{ou} \qquad
        x + 3 \geq 5  \implies x \geq 2
    \]
    \[
        \boxed{S = \left]-\infty,\ -8\right] \cup \left[2,\ +\infty\right[}
    \]

    \item $|2x - 1| > 3$
    \[
        2x - 1 < -3 \implies x < -1
        \qquad \text{ou} \qquad
        2x - 1 > 3  \implies x > 2
    \]
    \[
        \boxed{S = \left]-\infty,\ -1\right[ \cup \left]2,\ +\infty\right[}
    \]

    \item $|-x + 4| \leq 4$

    Note que $|-x+4| = |x-4|$, portanto:
    \[
        -4 \leq x - 4 \leq 4 \implies 0 \leq x \leq 8
    \]
    \[
        \boxed{S = \left[0,\ 8\right]}
    \]

\end{enumerate}
\pagebreak 

\section*{Exercício 7 --- Resolução}

\begin{enumerate}[label=(\alph*)]
    \item $p(x) = 3x + 2$: \hfill grau 1 (linear)
    \item $q(x) = x^2 - 4x + 1$ \hfill grau 2 (quadrático)
    \item $r(x) = 2x^3 - x + 5$ \hfill grau 3 (cúbico)
    \item $s(x) = 7$ \hfill grau 0 (constante)
    \item $p(x,y) = x^2 + y^2$ \hfill grau 2
    \item $q(x,y) = xy + x$ \hfill grau 2 \quad ($xy$: $1+1=2$)
    \item $r(x,y) = x^3 + y$ \hfill grau 3
    \item $s(x,y) = 2x^2y$ \hfill grau 3 \quad ($x^2y$: $2+1=3$)
\end{enumerate}

\section*{Exercício 8 --- Resolução}

\begin{enumerate}[label=(\alph*)]

    \item $2x^2 - 3x + 1 = 0$

    Forma padrão $ax^2 + bx + c$:
    \[
        a = 2, \quad b = -3, \quad c = 1
    \]

    \item $-x^3 + 4x = 7 \implies -x^3 + 4x - 7 = 0$

    Forma padrão $ax^3 + bx^2 + cx + d$:
    \[
        a = -1, \quad b = 0, \quad c = 4, \quad d = -7
    \]

    \item $5x^5 + 2x^2 = 0$

    Forma padrão $ax^5 + bx^4 + cx^3 + dx^2 + ex + f$:
    \[
        a = 5,\ b = 0,\ c = 0,\ d = 2,\ e = 0,\ f = 0
    \]

\end{enumerate}

\section*{Exercício 9 --- Resolução}

\begin{enumerate}[label=(\alph*)]

    \item $x^2 - 2x + 3 = 0$
    \[
        a = 1, \quad b = -2, \quad c = 3
    \]

    \item $3x^2 + x - 5 = 0$
    \[
        a = 3, \quad b = 1, \quad c = -5
    \]

    \item $-2x^2 + 4 = 0$
    \[
        a = -2, \quad b = 0, \quad c = 4
    \]

\end{enumerate}

\section*{Exercício 10 --- Resolução}

Utilizando a fórmula de Bhaskara: $x = \dfrac{-b \pm \sqrt{\Delta}}{2a}$, onde $\Delta = b^2 - 4ac$.

\begin{enumerate}[label=(\alph*)]

    \item $x^2 - 4 = 0$

    $a = 1,\ b = 0,\ c = -4$
    \[
        \Delta = 0^2 - 4(1)(-4) = 16
    \]
    \[
        x = \frac{0 \pm \sqrt{16}}{2} = \frac{\pm 4}{2}
    \]
    \[
        \boxed{x_1 = -2, \quad x_2 = 2}
    \]

    \item $x^2 - 9 = 0$

    $a = 1,\ b = 0,\ c = -9$
    \[
        \Delta = 0^2 - 4(1)(-9) = 36
    \]
    \[
        x = \frac{0 \pm \sqrt{36}}{2} = \frac{\pm 6}{2}
    \]
    \[
        \boxed{x_1 = -3, \quad x_2 = 3}
    \]

    \item $x^2 + 5x + 6 = 0$

    $a = 1,\ b = 5,\ c = 6$
    \[
        \Delta = 5^2 - 4(1)(6) = 25 - 24 = 1
    \]
    \[
        x = \frac{-5 \pm \sqrt{1}}{2} = \frac{-5 \pm 1}{2}
    \]
    \[
        x_1 = \frac{-5 - 1}{2} = \frac{-6}{2} = -3
        \qquad
        x_2 = \frac{-5 + 1}{2} = \frac{-4}{2} = -2
    \]
    \[
        \boxed{x_1 = -3, \quad x_2 = -2}
    \]

    \item $x^2 - x - 6 = 0$

    $a = 1,\ b = -1,\ c = -6$
    \[
        \Delta = (-1)^2 - 4(1)(-6) = 1 + 24 = 25
    \]
    \[
        x = \frac{1 \pm \sqrt{25}}{2} = \frac{1 \pm 5}{2}
    \]
    \[
        x_1 = \frac{1 - 5}{2} = \frac{-4}{2} = -2
        \qquad
        x_2 = \frac{1 + 5}{2} = \frac{6}{2} = 3
    \]
    \[
        \boxed{x_1 = -2, \quad x_2 = 3}
    \]

    \item $x^2 - 2x = 0$

    $a = 1,\ b = -2,\ c = 0$
    \[
        \Delta = (-2)^2 - 4(1)(0) = 4
    \]
    \[
        x = \frac{2 \pm \sqrt{4}}{2} = \frac{2 \pm 2}{2}
    \]
    \[
        x_1 = \frac{2 - 2}{2} = 0
        \qquad
        x_2 = \frac{2 + 2}{2} = 2
    \]
    \[
        \boxed{x_1 = 0, \quad x_2 = 2}
    \]

    \item $x^2 + 4x = 0$

    $a = 1,\ b = 4,\ c = 0$
    \[
        \Delta = 4^2 - 4(1)(0) = 16
    \]
    \[
        x = \frac{-4 \pm \sqrt{16}}{2} = \frac{-4 \pm 4}{2}
    \]
    \[
        x_1 = \frac{-4 - 4}{2} = \frac{-8}{2} = -4
        \qquad
        x_2 = \frac{-4 + 4}{2} = 0
    \]
    \[
        \boxed{x_1 = -4, \quad x_2 = 0}
    \]

\end{enumerate}
\end{document}
